%%==================================================
%% conclusion.tex for BIT Master Thesis
%% modified by yang yating
%% version: 0.1
%% last update: Dec 25th, 2016
%%==================================================
% !TeX root = ../main.tex
\chapter{总结与展望}
\enchapter{Conclusion and Future Work}
\label{chap:conclusion}

\section{工作总结}
\ensection{Conclusion}

随着人工智能、机器人技术以及虚拟仿真技术的快速发展，面向化学实验的具身智能研究正逐渐成为自动化实验、智能实验室和数字化科研中的重要方向。然而，现有相关研究在化学实验场景生成、高保真化学资产重建以及面向化学实验的具身智能任务规划等方面仍存在明显不足，限制了具身智能系统在化学实验任务中的训练效率、仿真真实性和实际应用能力。针对上述问题，本文围绕面向化学实验的具身智能关键技术，开展了化学实验场景生成、高保真化学资产重建以及面向化学实验的具身智能任务规划三方面研究，主要完成的工作如下：

（1）本文提出了基于大语言模型的化学实验场景生成算法。针对传统场景生成算法难以同时满足化学安全规范、场景质量与多样性的问题，本文构建了一种由大语言模型驱动的化学实验场景生成框架。该框架首先利用大语言模型解析自然语言形式的实验描述，生成实验物体清单与实验步骤，并结合化学知识库补全必要的安全设施。随后，系统根据实验流程与物体关系构建物体约束图，并结合资产检索与布局求解生成满足实验逻辑与安全要求的初始场景。考虑到复杂实验场景中人工自由浏览成本大、效率低，本文进一步研究了面向场景展示的相机规划算法，通过综合考虑关键区域显著性、场景覆盖率和相机运动平滑性，生成高效展示场景显著区域的相机路径。最后，为了进一步提升场景质量并增强训练环境的多样性，本文提出面向具身智能训练的场景进化算法，在安全约束基础上，结合大语言模型驱动的交叉与变异算子以及人工反馈，迭代生成兼顾质量与差异性的场景集合。该工作为具身智能体提供了更加安全、高质量的仿真实验环境。

（2）本文提出了基于隐式神经场的高保真化学资产重建算法。围绕化学资产重建这一目标，针对化学实验场景中透明物体和小物体难以高质量重建，以及神经隐式重建结果缺少显式纹理贴图、难以直接用于下游仿真平台的问题。本文构建了一条从多视角图像输入到显式纹理模型输出的完整重建流程。在数据预处理阶段，本文通过 COLMAP 恢复相机参数，并结合主体Mask完成稀疏点云去噪与空间归一化，为后续隐式场优化建立稳定的坐标参考。在透明物体重建方面，针对传统 NeuS 直接采用零等值面提取表面时容易出现表面缺失和内部伪影的问题，本文提出了基于绝对距离场局部极小值搜索的统一表面提取策略，实现了对透明区域与不透明区域的统一恢复。在小物体重建方面，针对目标像素占比较低、训练过程容易长时间停留于过渡阶段的问题，本文提出学习态增强策略，通过动态调节采样方向与距离场梯度方向之间的作用强度，加快训练初期的形状收敛，从而提高有限训练预算下的小物体几何重建质量。在纹理生成方面，本文进一步构建了从 NeuS 隐式颜色场到高分辨率纹理贴图的自动化生成算法，利用 Xatlas 完成 UV 参数化，并结合光栅化烘焙、形态学修复与 KNN 补全生成可直接用于仿真平台和渲染引擎的纹理资产。该工作为化学实验场景提供了高保真、仿真平台可用的三维资产基础。

（3）本文提出了从实验手册文本到机器人操作指令的分层转化算法。针对化学实验步骤通常以自然语言形式给出、难以直接转化为机器人可执行操作序列的问题，本文采用“实验手册文本—步骤计划—操作原语”的分层表示思路，实现从高层实验语义到机器人操作序列的逐级转化。首先，系统利用大语言模型对实验手册进行解析，提取实验过程中的关键操作、涉及对象及其目标要求，并生成结构化步骤计划；随后，基于预定义的机器人操作原语集合，将步骤计划进一步编译为仿真平台可调用的操作原语，从而建立从高层实验语义到机器人动作表示之间的映射。围绕化学实验中的典型操作，本文对抓取、放置、倾倒、搅拌四类基础实验操作原语进行了统一封装，使大语言模型能够在离散原语空间内完成实验步骤的组织与编排，而底层控制器则负责将原语接口转化为具体运动执行。最后，本文在基于 NVIDIA Omniverse 的机器人仿真平台中构建化学实验执行环境，并通过典型无机化学反应案例展示以及生成质量、生成效率的定量实验，对所提分层转化算法进行了系统验证。该工作为面向化学实验自动化的具身智能机器人系统提供了一种清晰、稳定且可扩展的任务规划框架。

\section{未来展望}
\ensection{Future Work}

尽管本文围绕化学实验场景生成、高保真化学资产重建以及机器人任务规划开展了较为系统的研究，并在虚拟仿真环境中取得了较好的实验结果，但整体上仍然存在进一步提升的空间。未来工作可从以下几个方向继续展开：

（1）面向化学实验全流程的多模态闭环场景生成。本文目前的场景生成主要以自然语言实验描述为核心输入，并通过人机协作的进化算法优化提升结果质量。未来可以结合多模态大模型，引入图像反馈，实现对生成场景的自动评估与迭代优化。同时，在仿真环境中引入机器人，构建机器人任务执行动态闭环优化机制，使生成场景不仅在静态空间关系上合理，而且能够在真实任务执行过程中持续根据可达性、碰撞风险、观测盲区以及任务成功率等反馈自动调整布局策略，从而进一步提升场景对具身智能训练的适配能力。

（2）面向复杂材质与高效率需求的高保真三维重建。本文所提出的重建算法在透明物体与小物体重建方面取得了一定效果，但对于具有更复杂光学性质的透明、半透明、高反光或混合材质物体，仍然存在进一步提升空间。未来可考虑在隐式表示中引入更强的物理先验或材质分解机制，将几何、反射、透射以及视角相关外观进行更细致的建模，以提升复杂化学器具的重建稳定性与外观真实性。此外，当前算法仍建立在较传统的 NeuS 训练框架之上，计算成本相对较高。后续可以结合 Instant-NGP 等高效神经表示算法，探索更加快速的训练与重建方案，并进一步研究隐式表示与显式表示相结合的混合建模算法，以减少从连续场到网格纹理模型转换过程中的信息损失，提升重建结果在仿真、渲染与机器人感知中的综合表现。

（3）面向真实自动化实验的闭环机器人执行与仿真迁移。本文在任务规划部分主要验证了从实验手册文本到操作原语的映射能力，并在虚拟仿真平台中展示了典型实验流程的执行效果。未来可在此基础上进一步研究结合环境感知的闭环执行机制，使机器人不仅能够按照预定操作原语完成操作，还能够根据视觉、力觉或状态反馈对实验过程中的异常情况进行检测与调整，例如容器偏移、液体未完全倒出、搅拌位置偏差或目标物体遮挡等问题。同时，后续还可探索从虚拟仿真平台到真实实验平台的迁移算法，结合数字孪生技术、领域随机化与参数校准策略，逐步缩小仿真与真实实验环境之间的差异，推动本文算法从虚拟验证走向真实自动化实验应用。此外，针对大语言模型在复杂化学任务规划中可能产生逻辑幻觉的问题，后续真实系统还需要建立独立于模型推理结果的物理级安全约束机制。该机制可结合化学规则硬约束、传感器状态监测、重量与液位核验、硬件急停和设备互锁等手段，形成不可被模型输出绕过的安全红线，从而提升系统在高危化学实验场景中的安全性与可靠性。进一步地，还可以研究多机器人协作、人机协同实验以及更复杂实验流程中的任务分解与执行问题，从而提升系统在实际化学实验场景中的适用范围。

